\documentclass[../../main.tex]{subfiles}
\graphicspath{{../../image/}{../../figure/}}
\begin{document}

\section{连续介质热力学的平衡律}\label{section:ch02-连续介质热力学平衡律}

连续介质热力学在本书中始终作为一个代表性模型，它由质量、线动量、角动量、能量和熵的平衡律所划定，下面将引入这些平衡律的参考与空间场方程。

在\emph{质量平衡律}中，既无通量也无产生，因此参考与空间场方程为
\begin{align}
\dot{\rho}_{0}=0,
\label{align:ch02-2-3-1}
\end{align}
\begin{align}
\rho_{t}+\mathrm{div}\,(\rho v^{\top})=0,
\label{align:ch02-2-3-2}
\end{align}
其中 $\rho_{0}$ 是\emph{参考密度}，$\rho$ 是与运动相伴随的\emph{密度}，二者由下式联系：
\begin{align}
\rho=\rho_{0}(\det F)^{-1}.
\label{align:ch02-2-3-3}
\end{align}
注意\eqref{align:ch02-2-3-1} 蕴涵与某质点相伴随的参考密度之值不随时间改变：$\rho_{0}=\rho_{0}(x)$。\eqref{align:ch02-2-3-2} 也称为\emph{连续性方程}。

在\emph{线动量平衡律}中，产生由（单位质量的）\emph{体力}向量 $b$（取值于 $\mathbb{R}^{m}$）诱导，而通量由取值于 $\mathbb{M}^{m\times m}$ 的\emph{应力张量}表示。参考与空间场方程为
\begin{align}
(\rho_{0}v)^{\cdot}=\mathrm{Div}\,S+\rho_{0}b,
\label{align:ch02-2-3-4}
\end{align}
\begin{align}
(\rho v)_{t}+\mathrm{div}\,(\rho vv^{\top})=\mathrm{div}\,T+\rho b,
\label{align:ch02-2-3-5}
\end{align}
其中 $S$ 表示\emph{Piola--Kirchhoff 应力}，$T$ 表示\emph{Cauchy 应力}，二者由下式联系：
\begin{align}
T=(\det F)^{-1}SF^{\top}.
\label{align:ch02-2-3-6}
\end{align}
对任一单位向量 $v$，$S v$ 在 $(x,t)$ 处的值给出在质点 $x$、时刻 $t$ 通过法向为 $v$ 的物质面传递的应力（单位面积的力）向量；而 $T v$ 在 $(\chi,t)$ 处的值给出在空间点 $\chi$、时刻 $t$ 通过法向为 $v$ 的空间面传递的应力向量。

在\emph{角动量平衡律}中，产生与通量分别是线动量平衡中所涉及的产生与通量关于原点的矩。因此，参考场方程为
\begin{align}
(\chi\wedge\rho_{0}v)^{\cdot}=\mathrm{Div}\,(\chi\wedge S)+\chi\wedge\rho_{0}b,
\label{align:ch02-2-3-7}
\end{align}
其中 $\wedge$ 表示外积。在 $\rho_{0}v$、$S$、$\rho_{0}b\in L^{1}_{\mathrm{loc}}$ 而运动 $\chi$ Lipschitz 连续的假设下，可利用\eqref{align:ch02-2-3-4}、\eqref{align:ch02-2-1-1} 和 \eqref{align:ch02-2-1-2} 将\eqref{align:ch02-2-3-7} 化为
\begin{align}
SF^{\top}=FS^{\top}.
\label{align:ch02-2-3-8}
\end{align}
类似地，角动量平衡的空间场方程由\eqref{align:ch02-2-3-5} 化为 Cauchy 应力张量对称这一陈述：
\begin{align}
T^{\top}=T.
\label{align:ch02-2-3-9}
\end{align}
无需进行该计算，因为\eqref{align:ch02-2-3-9} 也直接由\eqref{align:ch02-2-3-6} 和 \eqref{align:ch02-2-3-8} 推出。

在\emph{能量平衡律}中，能量密度是（单位质量的）\emph{内能} $\varepsilon$ 与动能之和。产生是体力功率与（单位质量的）\emph{热量供给} $r$ 之和。最后，通量是应力张量功率与\emph{热通量}之和。因此参考与空间场方程为
\begin{align}
(\rho_{0}\varepsilon+\tfrac{1}{2}\rho_{0}|v|^{2})^{\cdot}=\mathrm{Div}\,(v^{\top}S+Q^{\top})+\rho_{0}v^{\top}b+\rho_{0}r,
\label{align:ch02-2-3-10}
\end{align}
\begin{align}
(\rho\varepsilon+\tfrac{1}{2}\rho|v|^{2})_{t}+\mathrm{div}\,[(\rho\varepsilon+\tfrac{1}{2}\rho|v|^{2})v^{\top}]=\mathrm{div}\,(v^{\top}T+q^{\top})+\rho v^{\top}b+\rho r,
\label{align:ch02-2-3-11}
\end{align}
其中取值于 $\mathbb{R}^{m}$ 的参考与空间热通量向量 $Q$ 和 $q$ 由下式联系：
\begin{align}
q=(\det F)^{-1}FQ.
\label{align:ch02-2-3-12}
\end{align}

最后，\emph{熵平衡律}由 Clausius--Duhem 不等式
\begin{align}
(\rho_{0}s)^{\cdot}\geqslant\mathrm{Div}\,\left(\frac{1}{\theta}Q^{\top}\right)+\rho_{0}\frac{r}{\theta},
\label{align:ch02-2-3-13}
\end{align}
\begin{align}
(\rho s)_{t}+\mathrm{div}\,(\rho sv^{\top})\geqslant\mathrm{div}\,\left(\frac{1}{\theta}q^{\top}\right)+\rho\frac{r}{\theta},
\label{align:ch02-2-3-14}
\end{align}
分别以参考与空间形式表达。符号 $s$ 表示（单位质量的）\emph{熵}，$\theta$ 表示（绝对）\emph{温度}。于是\emph{熵通量}就是热通量除以温度。项 $r/\theta$ 表示由热量供给 $r$ 诱导的（单位质量的）外熵供给。然而，\eqref{align:ch02-2-3-13} 和 \eqref{align:ch02-2-3-14} 仅仅是不等式而非等式这一事实表明，可能存在额外的\emph{内熵产生}，它在本理论的框架中并未预先指定，只受非负性约束。这后一条件由热力学第二定律所规定（并且实际上表达了第二定律）。作为非负分布，内熵产生必然是一个测度 $\mathscr{N}$。将 $\mathscr{N}$ 加到右端，就把 Clausius--Duhem 不等式变成一个等式；由\cref{chapter:ch01-平衡律}，它是某平衡律的场方程。特别地，这表明即使当场仅仅局部可积时，参考形式\eqref{align:ch02-2-3-13} 与空间形式\eqref{align:ch02-2-3-14} 也是等价的。


运动与熵（或温度）场共同构成一个\emph{热力学过程}。内能、应力、热通量和温度（或熵）这些场由表征物体材料响应的本构关系从热力学过程中确定。特别地，应力的本构方程必须恒等地满足由\eqref{align:ch02-2-3-8} 或 \eqref{align:ch02-2-3-9} 所表达的角动量平衡律。有代表性的材料类别将在随后的\cref{section:ch02-热弹性} 和 \cref{section:ch02-热粘弹性} 中引入。

质量、线动量和能量平衡律的场方程与本构关系耦合，构成一个封闭的演化方程组，它应从给定的体力场 $b$、热量供给场 $r$、边界条件和初始条件出发确定热力学过程。

其余的熵平衡律起着截然不同的作用。Clausius--Duhem 不等式\eqref{align:ch02-2-3-13} 或 \eqref{align:ch02-2-3-14} 被视为对已经满足质量、动量和能量平衡律的热力学过程的\emph{热力学可容许性}判据。在这方面，光滑热力学过程与带间断的热力学过程被区别对待。

连续热力学的一个信条是：本构关系应当受到如下要求的约束，即任何平衡了质量、动量和能量的光滑热力学过程都必须自动是热力学可容许的。为实现这一要求，第一步是从 Clausius--Duhem 不等式导出\emph{耗散不等式}
\begin{align}
\rho_{0}\dot{\varepsilon}-\rho_{0}\theta\dot{s}-\mathrm{tr}\,(S\dot{F}^{\top})-\frac{1}{\theta}Q\cdot G\leqslant 0,
\label{align:ch02-2-3-15}
\end{align}
\begin{align}
\rho\dot{\varepsilon}-\rho\theta\dot{s}-\mathrm{tr}\,(TD)-\frac{1}{\theta}q\cdot g\leqslant 0,
\label{align:ch02-2-3-16}
\end{align}
分别为 Lagrange 或 Euler 形式，其中不涉及额外指定的体力和热量供给。出现在\eqref{align:ch02-2-3-15} 和 \eqref{align:ch02-2-3-16} 中的新符号 $G$ 和 $g$ 表示\emph{温度梯度}：
\begin{align}
G=\mathrm{Grad}\,\theta,\qquad g=\mathrm{grad}\,\theta,\qquad G=F^{\top}g.
\label{align:ch02-2-3-17}
\end{align}

为建立\eqref{align:ch02-2-3-15}，首先在质量、线动量和能量平衡律的场方程\eqref{align:ch02-2-3-1}、\eqref{align:ch02-2-3-4} 和 \eqref{align:ch02-2-3-10} 之间消去体力 $b$，得到
\begin{align}
\rho_{0}\dot{\varepsilon}=\mathrm{tr}\,(S\dot{F}^{\top})+\mathrm{Div}\,Q^{\top}+\rho_{0}r,
\label{align:ch02-2-3-18}
\end{align}
然后在上式与 Clausius--Duhem 不等式\eqref{align:ch02-2-3-13} 之间消去热量供给 $r$。类似地，\eqref{align:ch02-2-3-16} 由将\eqref{align:ch02-2-3-2}、\eqref{align:ch02-2-3-5}、\eqref{align:ch02-2-3-11} 与 \eqref{align:ch02-2-3-14} 相结合以消去 $b$ 和 $r$ 而得到。当然，\eqref{align:ch02-2-3-15} 与 \eqref{align:ch02-2-3-16} 等价：由\eqref{align:ch02-2-3-3}、\eqref{align:ch02-2-3-6}、\eqref{align:ch02-2-3-17}、\eqref{align:ch02-2-1-9} 和 \eqref{align:ch02-2-3-9}，任一者都蕴涵另一者。在上述计算中，至关重要的是假定所考虑的热力学过程是光滑的，因为这使我们能对 $|v|^{2}$、$v^{\top}S$、$\theta^{-1}Q$ 等项应用经典的乘积求导法则，从而产生大量抵消。应当强调，耗散不等式\eqref{align:ch02-2-3-15} 和 \eqref{align:ch02-2-3-16} 一般对带间断的热力学过程没有意义。

本构方程必须恒等地满足耗散不等式\eqref{align:ch02-2-3-15} 或 \eqref{align:ch02-2-3-16}，这将保证任何平衡了质量、动量和能量的光滑热力学过程自动是热力学可容许的。这一要求对具体材料类的实现将在随后的\cref{section:ch02-热弹性} 和 \cref{section:ch02-热粘弹性} 中展示。

除上述处理光滑热力学过程之外，Clausius--Duhem 不等式还负有鉴定间断过程热力学可容许性的额外责任。这是一个核心问题，具有许多方面，将在本书其余部分反复出现。

当处理具有复杂结构的连续介质时（例如不同材料的混合物），有必要用一个更一般的熵不等式代替 Clausius--Duhem 不等式，其中熵通量不再预先取为热通量除以温度，而是由一个单独的本构关系指定。然而结果表明，在热弹性或热粘弹性介质（本书主要关注的对象）的框架中，要求这样一个不等式对任何平衡了质量、动量和能量的光滑热力学过程恒成立，特别蕴涵熵通量必然是热通量除以温度，于是我们回到经典的 Clausius--Duhem 不等式。

为给即将进行的材料对称性研究做准备，有必要讨论当参考构型经历一个由等容双 Lipschitz 同胚 $\bar{x}$ 诱导的变化时，平衡律中所涉场的变换律，其 Jacobi 矩阵 $H$ 是幺模的\eqref{align:ch02-2-2-5}；见\cref{figure:ch02-等容变换示意图}。变形梯度 $F$ 和拉伸张量 $D$（参见\eqref{align:ch02-2-1-9}）将变换为新场 $\bar{F}$ 和 $\bar{D}$：
\begin{align}
\bar{F}=FH^{-1},\qquad\bar{D}=D.
\label{align:ch02-2-3-19}
\end{align}
平衡律所涉及的参考密度 $\rho_{0}$、内能 $\varepsilon$、Piola--Kirchhoff 应力 $S$、熵 $s$、温度 $\theta$、参考热通量向量 $Q$、密度 $\rho$、Cauchy 应力 $T$ 和空间热通量向量 $q$ 也将按照规则\eqref{align:ch02-2-2-6} 或 \eqref{align:ch02-2-2-7} 变换为新场 $\bar{\rho}_{0},\bar{\varepsilon},\bar{S},\bar{s},\bar{\theta},\bar{Q},\bar{\rho},\bar{T},\bar{q}$，即
\begin{align}
\bar{\rho}_{0}=\rho_{0},\quad\bar{\varepsilon}=\varepsilon,\quad\bar{S}=SH^{\top},\quad\bar{s}=s,\quad\bar{\theta}=\theta,\quad\bar{Q}=HQ,
\label{align:ch02-2-3-20}
\end{align}
\begin{align}
\bar{\rho}=\rho,\qquad\bar{T}=T,\qquad\bar{q}=q.
\label{align:ch02-2-3-21}
\end{align}
此外，参考与空间温度梯度 $G$ 和 $g$ 将变换为 $\bar{G}$ 和 $\bar{g}$，且
\begin{align}
\bar{G}=(H^{-1})^{\top}G,\qquad\bar{g}=g.
\label{align:ch02-2-3-22}
\end{align}

\end{document}
